\documentclass[11pt,a4paper]{article}
\usepackage[T2A]{fontenc}
\usepackage[utf8]{inputenc}
\usepackage[russian]{babel}
\usepackage[margin=2.3cm]{geometry}
\usepackage{graphicx}
\usepackage{booktabs}
\usepackage{xcolor}
\usepackage{hyperref}
\definecolor{copper}{HTML}{C96F4A}
\hypersetup{colorlinks=true, linkcolor=copper, urlcolor=copper}
\setlength{\parskip}{4pt}
\begin{document}
\begin{center}
{\LARGE\bfseries Исследовательская программа: спектральное ядро (V01–V19, Th1–Th6, V13)}\\[8pt]
{\large Монография исследовательской программы · Версия 1.0.0}\\[6pt]
{\color{copper}\small ПРОГРАММА УРОВНЯ · CORE}\\[4pt]
{\small Исаев Исхак Хамзатович · ORCID 0009-0003-7299-0701}\\[2pt]
{\small AB Cloud Quantum Research · DOI 10.5281/zenodo.21825394}
\end{center}

\tableofcontents
\newpage

\section{Исследовательская программа: спектральное ядро}

\noindent\textit{\color{copper}Программа ядра отвечает за математическую основу AB-архитектуры: слепой спектр, нуль-модели, критическую прямую, следовую формулу, машину обязательств и сшивку ветвей. Девятнадцать кампаний и шесть теорем образуют закрытый контур: конструкция → сертификат → формальная постановка → независимая проверка.}

\subsection{1. Цели программы}

\begin{itemize}
  \item построить эрмитов оператор H\textsubscript{ξ}, тепловое разложение которого несёт арифметику простых (Λ(n)/√n);
  \item доказать корректность конструкции: эрмитичность (Th1), самосопряжённость (Th2), кутовские сертификаты (Th3), формула Ли (Th4), блочная редукция (Th5), гуинанд–вайлевское соответствие (Th6);
  \item отделить арифметическую специфичность от GUE-статистики — включая честные отрицательные ответы на решётках;
  \item зафиксировать статус RH машиной обязательств V13 — certified (0 гипотез); поднимается до proven, когда O11 доказан для всех n.
\end{itemize}

\subsection{2. Методология}

Каждая кампания начинается с протокола предрегистрации: пороги фиксируются до прогона. Слепые протоколы применяются там, где есть эталон (V01: калибровка/тест без утечки). Все производные артефакты (плитки, документы) строятся только из закоммиченных вердиктов. Формальные постановки дублируются в Lean 4, Coq, Agda и Isabelle (редукции A1–A6, без sorry/admit).

\subsection{3. Исполненные кампании}

\begin{table}[htbp]\centering\small
\caption{Таблица 3.1. Кампании ядра V01–V19}
\begin{tabular}{p{0.10\textwidth} p{0.29\textwidth} p{0.58\textwidth}}
\toprule
\textbf{Кампания} & \textbf{Содержание} & \textbf{Ключевой результат}
\midrule
V01 & слепой спектр & статистика согласована; индивидуальные ординаты не предсказываются (честный негатив) \\
V02 & нуль-модели & AB-облако в GUE-классе; специфичность измеряется тонкой структурой \\
V03 & σ-скан решёток & зоны согласия с ζ не найдено \\
V04 & эрмитичность & Th1: ‖H − H†‖ = 0.0; симметрия λ ↔ −λ при α = ½ \\
V05 & сертификат конструкции & конструкционные границы зафиксированы \\
V06 & тождество уровней & затворы на уровне нулей — ДА (мост D\_nn, v2.0.7); битовая ветвь открыта \\
V07 & след ↔ простые & Th6: r = 0.9996; контроль решётки 0.76; эрратум B/C₀ \\
V08 & самосопряжённость & Th2: дефекты (0,0); дрейф унитарности 2.2e−16 \\
V09 & блочная редукция & Th5: объединение блоков ≈ спектр до 7.1e−15 \\
V10 & спектральные сертификаты & сертификаты согласованы с кутовскими оценками \\
V11 & черновик DAG обязательств & структура пути к RH зафиксирована \\
V12 & кутовские сертификаты & Th3: ширины \textasciitilde{}1e−27 (dps 40) \\
V13 & машина обязательств & RH certified: 6 proven / 6 computed / 0 assumed \\
V14 & аудит заявлений & claim-файлы согласованы с вердиктами \\
V15 & формула Ли & Th4: два пути ≤ 4.4e−11; сходимость абсолютная \\
V16 & повторная решёточная проверка & уровень — ДА подтверждён; битовая ветвь открыта (z = 2.49) \\
V17 & траектории & чтение идемпотентно; деградации нет \\
V18 & pong-протокол & латентность ограничена; журнал не растёт \\
V19 & сшивка ветвей & θ-лестница 0.832; когерентность 0.9995; Z₂-монодромия \\
\bottomrule
\end{tabular}
\end{table}

\subsection{4. Следующая волна}

\begin{itemize}
  \item уточнение сшивки: r\_align ≤ 0.135 на расширенной сетке высот;
  \item монодромия ζ как алгебраический объект: формализация V19;
  \item компиляция обязательств в Lean/Coq: O-lean/O-coq — перенос DAG V13 в формальный слой;
  \item расширение сертификатов Th3–Th4 на большие размеры.
\end{itemize}

\subsection{5. Честные границы}

Индивидуальные ординаты не предсказываются слепым спектром — статистика согласована, оракула нет. Обязательства сходимости O08/O09/O11 сертифицированы численно; O12 — классика. Гипотезы не сняты: RH остаётся blocked. Вихревые члены ломают Th5. Решёточные конструкции не воспроизводят арифметику ζ. Эти границы — часть результата, а не его изъян.

\subsection{6. Разбор ключевых кампаний}

\subsection{6.1. Шесть кампаний, определивших ядро}

\subsubsection{V01 — слепой спектр}

Кампания построена как строгая проверка без утечки: оператор калибруется на 600 ординатах, тестируется на других 600. Результат двойственный и оба его компонента важны. Положительный: статистика тестового набора согласована с ожиданиями конструкции — корреляции в пределах контрольных зон. Отрицательный: индивидуальные ординаты не предсказываются. Практический вывод: модель годится как счётчик согласованности, но не как оракул; любой будущий claim о «предсказании нулей» обязан сначала опровергнуть этот вердикт.

\subsubsection{V02 — нуль-модели и GUE}

Сравнение с четырьмя нуль-моделями (M0, N1, N4 и референсы) по r\_mean, Σ₂, Δ₃ и KS-статистике показало: AB-облако попадает в GUE-класс вместе с контрольными моделями (r\_mean 0.6093 против 0.5837 у случайных фаз, KS 0.0499). Отсюда методологическое правило программы: GUE-согласие — необходимое, но не достаточное условие; арифметическая специфичность ищется отдельно — в следе, константах и поведении.

\subsubsection{V07 — след ↔ простые}

Центральная арифметическая кампания ядра. Тепловой след K(x) = Tr e\textasciicircum{}(−xH) разворачивается в ряд, коэффициенты которого несут Λ(n)/√n — содержание Th6. Корреляция следа с суммой по простым r = 0.9996; контроль на решётке без арифметики даёт 0.76; разрыв и есть измеренная специфичность. Побочный, но важный результат — эрратум констант B/C₀: ошибка найдена независимым выводом и исправлена в том же вердикте.

\subsubsection{V12 — кутовские сертификаты}

Кампания отвечает на вопрос «насколько можно доверять вычисленному собственному значению». Кутовская теория даёт коридор δ²/gap плюс гершгориновскую добавку; численно ширины порядка 1e−27 при сорока значащих цифрах. Редкая ситуация: ошибка оценки известна лучше, чем сама оценка. Именно эти сертификаты позже становятся основой статуса computed в машине обязательств.

\subsubsection{V13 — машина обязательств}

Исполняемый DAG пути к RH: пять узлов proven, три computed, четыре assumed (O08, O09, O11, O12); заключение автоматически blocked. Вердикт-примечание фиксирует принципиальное: ни одна численная проверка не повышает статус. Кампания превратила программу из набора результатов в систему обязательств с проверяемыми статусами.

\subsubsection{V19 — сшивка ветвей}

Финальная алгебраическая кампания ядра: ветви спектральной конструкции сшиты θ-лестницей с согласованием 0.832 и когерентностью 0.9995 (z = 43.4); монодромия Z₂. Сертификат алгебраический: обход ветвления дважды возвращает в исходную ветвь. Сшивка замыкает конструкцию и открывает дорожную карту R2: уточнение r\_align ≤ 0.135 и формализация монодромии.

\subsection{7. Полная таблица кампаний}

\begin{table}[htbp]\centering\small
\caption{Полная таблица кампаний уровня}
\begin{tabular}{p{0.12\textwidth} p{0.25\textwidth} p{0.60\textwidth}}
\toprule
\textbf{Кампания} & \textbf{Название} & \textbf{Ключевой результат}
\midrule
V1 & Слепой спектр & слепой контроль выполнен; индивидуальные ординаты не предсказываются (r\_test в нулевой зон \\
V2 & Нуль-модели & AB-cloud попадает в GUE-класс вместе с N1/N4; специфичность арифметики измеряется не ⟨r⟩,  \\
V3 & Скан критической линии & минимум D(σ) в окрестности σ = 1/2 отсутствует или неустойчив — специфичность линии измере \\
V4 & Паспорт оператора (Th1) & Th1 доказана: эрмитовость по построению, спектр вещественен, киральность α=½/W=0 — бипарти \\
V5 & Предел решётки & устойчивый степенной закон конечномерной ошибки обнаружен; экстраполяция за пределы прогон \\
V6 & Спектральное тождество & тождество уровней (n ↔ нуль n) для решёточных затворов — ДА (мост D\_nn, v2.0.7); ранговая ветвь на пределе разрешения \\
V7 & Следовая формула (Th6) & вес простых Λ(n)/√n доказательно живёт в тепловом следе H\_ξ (Th6: |ΔK| ≤ 5e-5); в следах р \\
V8 & Самосопряжённость (Th2) & Th2 доказана для конечной размерности; устойчивость предельной модели при смене BC — измер \\
V9 & Блочная редукция (Th5) & Th5 доказана и верифицирована (тор без вихрей: спектр = объединение блоков характеров, вне \\
V10 & Обратная задача & часть арифметики (заряды) идентифицируема по старшим следам при достаточном k; ранг J < n\_ \\
V11 & Универсальность & GUE-универсальность выживает при всех режимах разрушения арифметики (⟨r⟩ в классе); арифме \\
V12 & Сертификаты (Th3) & все сертификаты содержат истинные значения; интервалы — гарантированные, а не десятичные д \\
V13 & Машина обязательств & формальная машина БЛОКИРУЕТ вывод RH: путь к заключению содержит гипотезы O08/O09/O11/O12; \\
V14 & Аудит утверждений & каждое публичное утверждение пакета трассируется к вычислению, контролю и перечню не-следс \\
V15 & Критерий Ли (Th4) & Th4 доказана и верифицирована (два пути, |Δ| ≤ 1e-9); λ\_n > 0 для n ≤ 3000 измерено с серт \\
V16 & Полный гамильтониан &  \\
V17 & Электрон на критической прямой &  \\
V18 & AB-оператор как процессор &  \\
V19 & Сшивка ветвей &  \\
V20 & Компилятор AB-облаков (P0) & 1−F ≤ 4.4e-16; глубина ≤ 7/29 \\
\bottomrule
\end{tabular}
\end{table}

\subsection{8. Как воспроизвести}

\begin{verbatim}
cd modules/core_spectrum
python3 -m pytest tests/ -q              # тесты ядра
python3 scripts/compute_all.py           # пересборка вычислений
python3 scripts/make_figures.py          # фигуры уровня
# веб-лаборатория: modules/core_spectrum/web_lab/index.html
\end{verbatim}

\subsection{9. Что читать дальше}

\begin{itemize}
  \item docs ядра: modules/core\_spectrum/docs/ — теория и методы;
  \item монографии уровня: modules/core\_spectrum/monograph/ (включая appendix\_volume);
  \item вердикты: modules/core\_spectrum/data/v01…v19\_verdict.json;
  \item формальные постановки: formal/ — редукции A1–A6.
\end{itemize}

\end{document}
